\begin{textsample}{POS Dim 1 – ro – Score 75.00 – ro/ro\_em\_10.txt}  \label{ex:f1_pos_131}
1.5.1.
O desenvolvimento dos jovens em suas múltiplas dimensões e a integração do currículo escolar em perspectiva à concepção de Educação Integral A Base Nacional Comum Curricular visa contemplar uma \textbf{abordagem} inicial ao repensar uma educação que não seja fragmentada, conteudista, disciplinar, que contemplará uma formação integral, visando transformar a vida dos estudantes, capacitando -os para lidar com os desafios e diferentes situações na \textbf{contemporaneidade}.
Posto que temos realidades em que pessoas são contratadas para trabalhar pelo seu conhecimento e logo são demitidas porque não têm atitudes suficientes de liderança em equipe, iniciativa de resolver situação-problema, entre outros.
Diante desse cenário, esse documento assume o compromisso com a educação integral, na perspectiva da busca do desenvolvimento humano global e na afirmação dos seus princípios : “ a superação da fragmentação radicalmente disciplinar do conhecimento, o estímulo à sua aplicação na vida real, protagonismo do \textbf{aluno} em sua aprendizagem e a importância do contexto para dar sentido ao que se aprende são alguns dos princípios subjacentes à BNCC. ” ( BRASIL, 2017, p. 17 ). \textbf{Posto} isso, a educação integral se \textbf{baseia} na concepção de um desenvolvimento pleno do ser humano, reconhece que só é possível essa integralidade quando os processos de aprendizagem ocorrem de modo multidimensional – física, afetiva, cognitiva, ética, estética e política – e se articulam com os diversos saberes da escola, da família, da comunidade e da região em que o indivíduo se insere.
Como concepção, a proposta de Educação Integral deve ser assumida por todos os agentes envolvidos no processo formativo das crianças, jovens e adultos.
Nesse contexto, a escola se converte em um espaço essencial para assegurar que todos os estudantes tenham uma formação integral garantida.
Ela assume o papel de articuladora das diversas experiências educativas em que os \textbf{alunos} podem viver dentro e fora dela, a partir de uma intencionalidade clara que favoreça as aprendizagens importantes para o seu desenvolvimento integral.
Assim, a fim de que essa formação ocorra, é necessário ser pautada conforme alguns princípios, tais como : centralidade do estudante, aprendizagem permanente, perspectiva inclusiva e gestão democrática.
O primeiro princípio é a centralidade no estudante.
O projeto pedagógico deve ser construído e revisitado, como também a proposta deve ser personalizada em suas necessidades e ter, de fato, a participação dos estudantes na construção de um processo de \textbf{ensino-aprendizagem} global.
Outro princípio é a aprendizagem permanente, o que pressupõe que todas as dimensões do processo de \textbf{ensino-aprendizagem} estejam inseridas no currículo.
Dessa forma, as dimensões desenvolvidas não devem ser somente a intelectual, mas também a social, emocional, física e cultural, compondo assim um desenvolvimento integral.
A perspectiva inclusiva também está entre os princípios da educação integral.
Portanto, as propostas pedagógicas devem respeitar todas as diferenças, como as deficiências, a origem étnica e racial, religiosa, entre outros.
Propõe -se que todos os espaços escolares sejam inclusivos e que, nesses locais, os estudantes tenham oportunidade de desenvolvimento em suas \textbf{inúmeras} dimensões.
Nesse sentido, é importante que haja uma gestão democrática, com o \textbf{intuito} de garantir os interesses e necessidades de aprendizagem e desenvolvimento dos estudantes.
Assim, a gestão democrática pressupõe que as decisões e o acompanhamento das atividades sejam realizados de forma coletiva com a comunidade escolar – \textbf{alunos}, pais e educadores.
Compreender a concepção de Educação em tempo integral, na perspectiva da Educação Integral, pressupõe o entendimento da própria visão de currículo e a formação do sujeito em sua plenitude, ou seja, omnilateral.
Nessa perspectiva, podemos referenciar a \textbf{compreensão} do homem omnilateral em Marx ( 2008 ), quando este destaca o fato de o homem usar todas as suas potencialidades no processo de transformação do meio para sua sobrevivência.
Para sobreviver, segundo Marx ( 2008 ), o homem lança mão de braços, pernas, \textbf{mente} e espírito, transforma a natureza humanizando -a e humanizando -se.
A omnilateralidade, portanto, não é o desenvolvimento de potencialidades humanas inatas.
É a criação dessas potencialidades pelo próprio homem, na prática social, no mundo do trabalho.
Nessa perspectiva, a educação está vinculada à produção social total, bem como a todas as dimensões da existência e da ação humana.
Seria conceber a plenitude humana, a \textbf{totalidade} de suas \textbf{capacidades} físicas, intelectuais e espirituais, dispondo do tempo de produção de fruição de cultura, concebendo essa como \textbf{inerente} às relações de trabalho.
Não se trata, portanto, de oferecer um tempo de \textbf{disciplinas} acadêmicas ou formais e outro de oficinas e atividades lúdicas, que geralmente se consubstanciam em trabalhos manuais e artísticos.
Por isso, essa formação deve ser estabelecida a partir de quatro pilares fundamentais a uma concepção de educação integral : \textbf{autoria}, ampliação de repertório, cocriação e autoformação.
Por \textbf{intermédio} desses princípios, busca -se estabelecer um diálogo com os jovens/estudantes ao direcionar o trabalho pedagógico a partir de ações em torno das mais diferentes expressões culturais, na perspectiva de valorizar a cultura juvenil dentro da comunidade escolar.
No entanto, há que se ter o cuidado na sua implementação, pois tais ações não podem assumir direções e alcances diferenciados.
Na escola, as práticas não podem ser reduzidas a um determinado tempo e espaço, como, por exemplo, em atividades extraescolares, fazendo delas um meio de ocupar o tempo dos estudantes, constituindo -se em um complemento, sem nenhum impacto no currículo.
Por \textbf{conseguinte}, nesse documento serão apresentadas as dez \textbf{competências} gerais estabelecidas na Base, que se \textbf{inter-relacionam} e perpassam todos os componentes curriculares para a construção de conhecimentos, habilidades, atitudes que, ao concluir a Educação Básica, deverão ser internalizadas pelos estudantes, conforme dispõe a BNCC que “ ser competente significa ser capaz de, ao se defrontar com um \textbf{problema}, ativar e utilizar o conhecimento construído ” ( BRASIL, 2017, p. 16 ).
É de fundamental importância \textbf{explicitar} que as \textbf{Competências} Gerais foram definidas a partir dos direitos éticos, estéticos e políticos assegurados pelas Diretrizes Curriculares Nacionais e dos conhecimentos, habilidades, atitudes e valores essenciais para a vida no \textbf{século} XXI.
Nesse sentido, o conjunto das dez \textbf{competências}, indicadas na BNCC, afirma a concepção de Educação Integral e a importância do conhecimento em uso, no lugar do acúmulo de informações, pois “ esse conjunto de \textbf{competências} \textbf{explicita} o compromisso da educação brasileira com a formação humana integral e com a construção de uma sociedade justa, democrática e inclusiva ” ( BRASIL, 2017, p. 19 ).
Com essas \textbf{competências}, a Base retoma uma perspectiva de educação para a vida que requer como princípios pedagógicos uma visão integrada da aprendizagem e do desenvolvimento humano, a \textbf{contextualização} na construção do conhecimento escolar e a participação ativa dos sujeitos envolvidos nos processos educativos, a começar pelos estudantes.
Por isso, ainda que as \textbf{Competências} Gerais sejam válidas para toda a Educação Básica, a sua \textbf{progressão} para o Ensino Médio deverá ser contemplada, pois no \textbf{decorrer} da Educação Infantil ao Ensino Médio, as \textbf{competências} permanecem iguais, mas se desdobram ao longo de cada uma dessas etapas da educação em diferentes direitos de aprendizagem, campos de experiência, unidades temáticas, objetos de conhecimento e habilidades específicas, adequando -se às particularidades de cada fase do desenvolvimento dos estudantes. \textbf{Posto} isso, \textbf{segue} o organograma abaixo que apresenta as dez \textbf{competências} gerais individualizadas por cores e enumeradas, \textbf{sinalizando} seus respectivos objetos. \textbf{Figura} 1 : \textbf{Competências} Gerais da BNCC Fonte : http://portal.educacao.rs.gov.br/Portals/1/Images/Novo\%20Ensino\%20M\%C3\%A9dio/2278.png Em seguida, apresentamos um \textbf{quadro} de \textbf{Competências} Gerais da Base Nacional Comum Curricular.
No Quadro 1 – \textbf{Competências} gerais estabelecidas pela Base – visam \textbf{detalhar} o objeto, o objetivo e a finalidade que compõem cada uma das 10 \textbf{Competências} Gerais da BNCC. \textbf{Quadro} 1 : \textbf{Competências} Gerais estabelecidas pela Base¹ - C1 – OBJETO DA \textbf{COMPETÊNCIA} : Conhecimento Objetivo : Valorizar e utilizar os conhecimentos sobre o mundo físico, social, cultural e \textbf{digital}.
Finalidade : Entender e explicar a realidade, continuar aprendendo e colaborar com a sociedade. - C2 – OBJETO DA \textbf{COMPETÊNCIA} : Pensamento científico, \textbf{crítico} e criativo.
Objetivo : Exercitar a curiosidade intelectual e utilizar as \textbf{ciências} com criticidade e criatividade.
Finalidade : Investigar causas, elaborar e testar hipóteses, formular e resolver \textbf{problemas} e criar soluções. - C3 – OBJETO DA \textbf{COMPETÊNCIA} : Responsabilidade e Cidadania.
Objetivo : Valorizar as diversas manifestações artísticas e culturais.
Finalidade : Fruir e participar de práticas diversificadas da produção \textbf{artístico-cultural}. - C4 – OBJETO DA \textbf{COMPETÊNCIA} : Comunicação Objetivo : Utilizar diferentes linguagens.
Finalidade : Expressar -se e partilhar informações, experiências, ideias, sentimentos e produzir sentidos que levem ao entendimento mútuo. - C5 – OBJETO DA \textbf{COMPETÊNCIA} : Cultura Digital Objetivo : Compreender, utilizar e criar \textbf{tecnologias} \textbf{digitais} de forma \textbf{crítica}, significativa e ética.
Finalidade : Comunicar -se, acessar e produzir informações e conhecimentos, resolver \textbf{problemas} e exercer protagonismo e \textbf{autoria}. - C6 – OBJETO DA \textbf{COMPETÊNCIA} : Trabalho e Projeto de Vida Objetivo : Valorizar e apropriar -se de conhecimentos e experiências.
Finalidade : Entender o mundo do trabalho e fazer escolhas alinhadas à cidadania e ao seu projeto de vida com liberdade, autonomia, criticidade e responsabilidade. - C7 – OBJETO DA \textbf{COMPETÊNCIA} : Argumentação Objetivo : \textbf{Argumentar} com base em fatos, dados e informações confiáveis.
Finalidade : Formular, negociar e defender ideias, pontos de \textbf{vista} e decisões comuns, com base em direitos humanos, consciência \textbf{socioambiental}, consumo responsável e ética. - C8 – OBJETO DA \textbf{COMPETÊNCIA} : Autoconhecimento e Autocuidado.
Objetivo : Conhecer -se, compreender -se na diversidade humana e apreciar -se.
Finalidade : Cuidar de sua saúde física e emocional, reconhecendo suas emoções e as dos outros, com autocrítica e \textbf{capacidade} para lidar com elas. - C9 – OBJETO DA \textbf{COMPETÊNCIA} : Empatia e Cooperação.
Objetivo : Exercitar a empatia, o diálogo, a resolução de conflitos e a cooperação.
Finalidade : Fazer -se respeitar e promover o respeito ao outro e aos direitos humanos, com acolhimento e valorização da diversidade, sem preconceitos de qualquer natureza. - C10 – OBJETO DA \textbf{COMPETÊNCIA} : Responsabilidade e Cidadania.
Objetivo : Agir pessoal e coletivamente com autonomia, responsabilidade, flexibilidade, resiliência e determinação.
Finalidade : Tomar decisões com base em princípios éticos, democráticos, inclusivos, sustentáveis e solidários. ¹ A letra C equivale à \textbf{Competência}.
Ao realizar a \textbf{análise} do \textbf{quadro}, é \textbf{imprescindível} destacar que as \textbf{competências} gerais, apresentadas no Quadro 1, \textbf{inter-relacionam} -se e desdobram -se no tratamento didático proposto para as três etapas da Educação Básica ( Educação Infantil, Ensino Fundamental e Ensino Médio ), articulando -se na construção de conhecimentos, no desenvolvimento de habilidades e na formação de atitudes e valores, nos termos da Lei de Diretrizes e Bases da Educação.
Na perspectiva de enriquecer a \textbf{análise} do \textbf{desdobramento} da temática em evidência, \textbf{sugere} -se que o Quadro 1 – \textbf{Competências} gerais estabelecidas pela Base – seja analisado com o olhar voltado para o componente curricular, buscando suas possibilidades na contribuição da formação integral a partir do ensino por \textbf{competências} conforme a BNCC conceitua : “ \textbf{competência} é definida como a \textbf{mobilização} de conhecimentos ( conceitos e procedimentos ), habilidades ( práticas, cognitivas e \textbf{socioemocionais} ), atitudes e valores para resolver \textbf{demandas} \textbf{complexas} da vida cotidiana, do pleno exercício da cidadania e do mundo do trabalho. ” ( BRASIL, 2017, p. 8 ).
Vale \textbf{salientar} também que a BNCC, ao definir as dez \textbf{competências}, reconhece a importância dos valores e o estímulo de ações que contribuam para a transformação da sociedade, tornando -a mais humana, justa e direcionada para a preservação da natureza.
A BNCC, como documento norteador, indica que as decisões pedagógicas devem estar orientadas para o desenvolvimento de \textbf{competências}.
Isso deve acontecer por meio da \textbf{indicação} clara do que os estudantes devem “ saber ” ( considerando a constituição de conhecimentos, habilidades, atitudes e valores ) e, sobretudo, do que devem “ saber fazer ” ( considerando a \textbf{mobilização} desses conhecimentos, habilidades, atitudes e valores para resolver \textbf{demandas} \textbf{complexas} da vida cotidiana, do pleno exercício da cidadania e do mundo do trabalho ).
Dessa forma, a explicitação das \textbf{competências} oferece referências para o fortalecimento de ações que assegurem as aprendizagens essenciais definidas na BNCC.
Ao postular \textbf{competências} gerais como objetivos \textbf{transversais} e integradores para o ensino e aprendizagem da construção de conhecimentos \textbf{relevantes} para a vida, a proposta curricular do estado de Rondônia também se dedicou, na sua construção, ao enfoque nas estratégias de aprendizagem ativa ( no uso e para o uso dos conhecimentos ) e nas \textbf{capacidades} de transferência de saberes de um contexto a outro, de um componente a outro, da sala de \textbf{aula} e da escola aos diferentes contextos sociais.
Portanto, a formação dos estudantes, orientada pelo desenvolvimento das \textbf{competências} gerais, pressupõe uma reflexão permanente sobre o currículo, principalmente em relação ao porquê, o quê e o como se ensinar.
Nesse sentido, as \textbf{competências} funcionam como macro-objetivos e condicionantes \textbf{metodológicos} das propostas integradoras de todo desenho curricular e da implementação das políticas educacionais correlacionadas : aprender conhecimentos, habilidades, valores e atitudes no uso e para o uso social dessas aprendizagens.
Aprender para \textbf{mobilizar} e transferir conhecimentos de um campo a outro do conhecimento e entre diferentes contextos, da escola às situações sociais cotidianas.
Habilitar crianças e jovens a serem \textbf{autores} e autônomos no uso das \textbf{tecnologias} para desenvolvimento da \textbf{compreensão}, \textbf{análise} e transferência de conhecimentos e \textbf{capacidades} a novas situações.
Aprender no contexto de uso e para o uso social de conhecimentos e habilidades, ou seja, estimular a criação e o compartilhamento de saberes e \textbf{capacidades} para além do contexto de aprendizagem inicial, aplicando -os em situações novas e contextualizadas.
O ensino por \textbf{competências} na Educação Integral não pode ser jamais a aplicação mecânica de metodologias que promovam habilidades em sala de \textbf{aula}.
Trata -se, sempre, de interpretar e refletir institucionalmente sobre os propósitos, o contexto e as formas dessa implementação como garantias de direitos de aprendizagem e desenvolvimento a serviço de um projeto formativo e ético para o desenvolvimento integral de todos e de cada um.
Para \textbf{tanto}, a formação na Educação Integral constrói -se como uma espécie de elo entre a concepção, o currículo e a avaliação.
Assim como a concepção curricular, a proposta de formação também convoca mudanças e estrutura condizentes com a visão de educação integral.
Nesse viés, o ambiente manifesta a intenção de educação humanizada, potencializadora da criatividade, disponibilizando os recursos para exploração, promovendo a convivência enriquecedora das diferenças, a apropriação dos diversos lugares de aprender do território e a relação sustentável com os recursos do planeta.
Nessa perspectiva, é preciso desenvolver estratégias que estimulem o diálogo entre os diversos segmentos da comunidade, a mediação de conflitos por pares, o bem-estar de todos, a valorização da diversidade e das diferenças, a promoção da equidade, da própria sustentabilidade e a integração com a comunidade.
Assim, não apenas o conteúdo de formação se adequa, mas a própria forma como se oferta é transformada.
Na concepção de educação para o desenvolvimento integral, o planejamento das atividades de ensino e aprendizagem deve valorizar a \textbf{autoria} dos professores, com o máximo de escuta e diálogo possível com seus pares e estudantes, com o objetivo de estimular ciclo de aprendizagem ativa e autônoma dos estudantes.
A experimentação, a personalização e a avaliação formativa são elementos importantes quando se quer efetivamente garantir propostas \textbf{contemporâneas} equitativas e inclusivas dentro e fora da sala de \textbf{aula} a partir de metodologias ativas.
Dessa forma, para trabalhar o desenvolvimento integral dos estudantes, a escola pode optar por disponibilizar o conhecimento com base nas áreas previstas nas Diretrizes Curriculares Nacionais para a Educação Básica, como, por exemplo, organizando \textbf{sequências} didáticas, projetos, roteiros de estudos ou pesquisas em cada área de conhecimento, mas buscando sempre a integração entre os diferentes componentes curriculares e garantindo que os estudantes possam também escolher temas de seu interesse.
Na perspectiva do progressivo desenvolvimento integral de \textbf{competências}, o modelo instrucional da \textbf{aula} expositiva tem seu papel bastante diminuído e passa a ser decisivo saber mediar e conduzir boas situações de comunicação, cooperação, colaboração e cocriação entre crianças e jovens que pertencem a grupos diversificados com graus de habilidade distintos.
Isto significa \textbf{dizer} que o currículo não pode ser concebido como um rol de \textbf{disciplinas} organizadas de forma linear no tempo já culturalmente definido em função de horas-aula.
Por outro \textbf{lado}, ele não pode ser tomado como se fosse um objeto per se.
O currículo não tem vida própria, ele não é a expressão dele mesmo.
Ele revela, então, um projeto de realidade, de mundo, de homem e de educação.
Nele se manifesta a concepção que se define a partir da intencionalidade histórica e social sobre a escola.
O currículo integrado, por sua vez, se caracteriza exatamente pela integração de todas as potencialidades da condição humana.
Assim, um currículo, na perspectiva da Educação Integral, deve integrar os \textbf{potenciais} educativos, pois isto amplia as ferramentas de \textbf{contextualização} no processo de produção do conhecimento, aumenta os vínculos estabelecidos entre os estudantes, diversifica as ofertas educativas e os espaços de aprendizagem, potencializando, entre outras coisas, a tão desejada articulação entre escola e famílias, escola e comunidade.
Pode -se afirmar, igualmente, que o currículo integrador também apoia o desenvolvimento das \textbf{competências} gerais da BNCC, na medida em que oferece um campo concreto de experiências no qual os estudantes podem aplicar e consolidar habilidades.
A fim de que a proposta do currículo integrado vigore, é importante que a investigação seja a característica predominante do docente a partir de uma prática pedagógica centralizada no protagonismo da aprendizagem dos estudantes, vinculada a um processo coletivo de acordos, pesquisas e efetivação de boas práticas.
Dessa forma, o trabalho pedagógico ganhará uma dimensão criadora totalmente estruturada no coletivo, levando em consideração os saberes já adquiridos dos estudantes e os questionamentos que os conduzem à pesquisa, garantindo a qualidade e a eficácia do sistema educacional em sua plenitude.
Sendo assim, a gestão de práticas colaborativas, ao se \textbf{materializarem} a partir dos princípios norteadores da Educação Integral, será \textbf{basilar} para um aprendizado significativo ao instaurar no cotidiano dos docentes um processo de reflexão contínua.
Referente ao Ensino Médio, o currículo será \textbf{permeado} pelas contribuições das comunidades escolares e, principalmente, dos estudantes, pois, como afirmam as Diretrizes Nacionais, “ Ninguém mais do que os participantes da atividade escolar em seus diferentes segmentos, conhecem a sua realidade, portanto, estão mais habilitados para tomar decisões a respeito do currículo que leva à prática. ” ( DCNEM, 2013, p. 183 ).
Assim, as novas formas de organização dos componentes curriculares e o reconhecimento da participação dos estudantes na escolha dos conteúdos são inovações \textbf{pertinentes} a um novo processo democratizante. \textbf{Posto} isso, o contexto de cada instituição escolar será \textbf{permeado} a partir de sentidos construídos, praticados e compreendidos de conhecimentos significativos de cada área do conhecimento ( linguagens, matemática, \textbf{ciências} humanas etc. ), resultando em um currículo integrado no qual a participação de todos os envolvidos é efetivada.
É \textbf{imprescindível} a superação de currículos engessados que projetam conteúdos com rigidez por meio de articulações sistemáticas sem a preocupação de atender aos anseios dos estudantes.
A intenção é que ocorra a articulação entre os componentes curriculares, permitindo a organização dos saberes e o favorecimento da diversidade com o \textbf{intuito} de produzir um elo e o diálogo entre todas as áreas do conhecimento.
A integração do trabalho pedagógico deve ser estruturada por meio da troca, da pesquisa e da \textbf{sistematização} dos saberes entre docentes e estudantes, com momentos de compartilhamento das experiências, de modo que todos os envolvidos se considerem sujeitos de sua aprendizagem.
Assim, o currículo escolar será considerado como um espaço de escolhas, experiências, por meio da seleção de saberes com o \textbf{escopo} de possibilitar a transição de conhecimentos mais profundos e abrangentes para a \textbf{compreensão} do meio em que está inserido.
Logo, o currículo integrado, ao abordar o protagonismo do conhecimento e transformar os objetos de estudo em necessidades dos jovens, proporcionará a sua formação integral.
Nessa perspectiva, a educação integral, para que se concretize, não deve permitir a ruptura entre a \textbf{capacidade} de criar com as relações de afetividade, pois o desenvolvimento pleno ocorre a partir do elo entre o jovem com ele mesmo e com o mundo.
São as sensações, os sentimentos e a criatividade que irão \textbf{agregar} as dimensões física, afetiva, intelectual e social ao desenvolvimento pleno do estudante.
A contribuição para efetivar o currículo integrado sustentar -se-á pela heterogeneidade das concepções de mundo, conforme afirma Paulo Freire ( 2005 ) ao relatar que o dialogismo na educação é representado pelas relações entre educador e educando a partir da mediatização do mundo, ou seja, esse “ mundo ” \textbf{configura} -se como objeto de pesquisa para ambos.
Assim, essa proposta proporcionará às escolas o acolhimento de diversidades locais dos envolvidos no processo de ensino e aprendizagem.
Concomitantemente, estimulará a valorização regional e o sentimento de pertencimento para uma \textbf{juventude} que tem a oportunidade de escolhas por meio de experiências, visões de mundo diferenciadas, ao conduzi -la a uma reflexão acerca dos seus relacionamentos com o \textbf{intuito} de \textbf{auxiliar} na construção de novos saberes.
Outro aspecto a ressaltar está relacionado à intersetorialidade, uma vez que o currículo deve interligar os estudantes a um espaço para a realização de diversificados projetos interdisciplinares, possibilitando a exploração, a pesquisa de distintos temas conforme os seus interesses, bem como a realização de projetos em \textbf{localidades} que circundam o seu território.
Dentre eles, se destacam centros de cultura e esporte, bibliotecas, museus, universidades.
Esse conceito é abordado nas Diretrizes Curriculares Nacionais ( DCNEM ) ao afirmar que o currículo está intrinsecamente ligado à intersetorialidade, tendo em \textbf{vista} que não há cisão entre o educar e o zelar, visto que : > O desenvolvimento da \textbf{abordagem} territorial pode levar à elaboração de planos educativos locais, \textbf{baseados} em \textbf{diagnósticos} coletivos sobre a realidade daquele espaço na definição de prioridades e no compartilhamento de responsabilidades, visando a constituição de condições para o desenvolvimento integral de todos os estudantes. ( DCNEM, 2013, p. 18 ) A transformação do território deve estar articulada aos saberes locais e às possibilidades educativas, sempre com o olhar atento no suporte da aprendizagem dos estudantes, a fim de que o território se torne um lugar \textbf{propício} para atos educativos.
Essa intencionalidade pedagógica oportunizará a construção de práticas fundamentadas em aproximações, \textbf{aceitações}, diversidades de linguagens e convívios significativos.
A constituição dessas práticas pode ser efetivada com a valorização dos elementos articuladores do currículo e território, dentre eles, se destaca o reconhecimento dos saberes locais, pois \textbf{auxiliam} na percepção de práticas do dia a dia pertencentes àquele território, tais como : valores, memórias, hábitos.
Diante dessa perspectiva, percebe -se quão \textbf{imprescindíveis} são os saberes locais na construção de aprendizagens significativas e \textbf{relevantes} para os jovens, tendo em \textbf{vista} que possibilitam a execução de práticas e aquisição de conhecimentos a partir das vivências dos estudantes.
É mediante a esses saberes que podem ser criados temas geradores a um contexto de reflexão referente à linguagem, à matemática, às \textbf{ciências} humanas e às \textbf{ciências} da natureza.
Para que os saberes locais sejam reconhecidos, é necessário que a intencionalidade pedagógica seja articulada aos objetivos do currículo ( \textbf{competências} gerais, conhecimentos e habilidades interdisciplinares ou propostos pelas áreas do conhecimento ).
No currículo integrado, agentes, espaços e dinâmicas de cada território devem ser considerados como \textbf{potenciais} educativos capazes de proporcionar aprendizagem, haja \textbf{vista} que eles evidenciam a concretização dos saberes locais e dos objetos de reflexão das áreas do conhecimento no território vivenciado.
São compreendidos como agentes pessoas, coletivos ou instituições pertencentes ao local que atuam direta ou indiretamente ao realizar alterações nas suas dinâmicas a partir de intervenções.
Por meio da articulação com esses agentes há a possibilidade da remodelagem desse território, como também a formação integral dos estudantes.
Diante desse contexto, torna -se necessária a realização frequente entre os agentes e comunidade escolar por meio de \textbf{diagnóstico} e de escuta, os quais possibilitarão o elo com os eixos temáticos, os temas geradores ou temas estruturantes, conforme a BNCC preconiza.
Tal \textbf{abordagem}, direcionada às práticas de ensino e aprendizagem, confere pertinência e relevância para contextualizar os conteúdos curriculares ( unidades temáticas, objetos de conhecimento e habilidades ) no território – mediante temas, questões, saberes e práticas significativas para os estudantes.

% matched lemmas: abordagem, aceitação, agregar, aluno, análise, argumentar, artístico-cultural, aula, autor, autoria, auxiliar, basear, basilar, capacidade, ciência, competência, complexo, compreensão, configurar, conseguinte, contemporaneidade, contemporâneo, contextualização, crítico, decorrer, demanda, desdobramento, detalhar, diagnóstico, digital, disciplina, dizer, ensino-aprendizagem, escopo, explicitar, figurar, imprescindível, indicação, inerente, inter-relacionar, intermédio, intuito, inúmero, juventude, lado, localidade, materializar, mente, metodológico, mobilizar, mobilização, permear, pertinente, postar, potencial, problema, progressão, propício, quadro, relevante, salientar, segar, sequência, sinalizar, sistematização, socioambiental, socioemocional, sugerir, século, tanto, tecnologia, totalidade, transversal, vista
\end{textsample}
